% !PRINTED 354
\weberpage{354}
%<<CONT
anderen Punkte als die Berührungspunkte der Doppeltangenten
befriedigt ist.

Die Variablen $y$ genügen der Tangentengleichung

\weq{7}{y_1 f'(x_1) + y_2 f'(x_2) + y_3 f'(x_3) = 0.}

Bezeichnen wir mit $b_1, b_2, b_3$ irgend drei willkürliche Constanten, und setzen

\weq{8}{b_x = b_1 x_1 + b_2 x_2 + b_3 x_3,}

so dass $b_x = 0$ die Gleichung einer willkürlichen geraden Linie ist, so können wir zu (7) noch die Gleichung

\weq{9}{b_1 y_1 + b_2 y_2 + b_3 y_3 = 0}

hinzunehmen, also $(y)$ als den Schnittpunkt der Tangente in $(x)$ mit der beliebigen geraden Linie $b_x$ auffassen. Dann erhalten wir

\weq{10}{\begin{aligned}
y_1 &= b_2 f'(x_3) - b_3 f'(x_2),\\
y_2 &= b_3 f'(x_1) - b_1 f'(x_3),\\
y_3 &= b_1 f'(x_2) - b_2 f'(x_1),
\end{aligned}}

und hierdurch ist die Bedingung (7) identisch befriedigt. Durch diese Substitution geht $R(x,y)$ in eine homogene Function $20^{\text{sten}}$ Grades der $x$ und $6^{\text{ten}}$ Grades der $b$ über, die wir mit $D(x,b)$ bezeichnen wollen. Diese Function $D(x,b)$ verschwindet aber ausser in den Berührungspunkten der Doppeltangenten noch in den vier Schnittpunkten der Geraden $b_x$ mit der Curve $f$. Diese Bedingung muss nun noch so umgeformt werden, dass sie von den $b$ unabhängig wird.

Gehen wir von dem Punkte $(y)$ zu einem beliebigen anderen Punkte der Tangente über, so können wir dies dadurch erreichen, dass wir $y$ durch $\mu x + \lambda y$ ersetzen, worin $\lambda, \mu$ zwei willkürliche Parameter bedeuten. Dadurch geht $f(x+ty)$ in

\[
f[(1+\mu t)x + \lambda t y] = (1+\mu t)^4 f\left(x + \frac{\lambda t}{1+\mu t}\, y\right)
\]

über, und die Entwickelung (3) ergiebt, wenn wir

\[
P_0(x,y),\ P_1(x,y),\ P_1(x,\mu x + \lambda y)
\]

gleich Null setzen:

\[
\begin{aligned}
&6t^2(1+\mu t)^2\lambda^2 P_2(x,y) + 4t^3(1+\mu t)\lambda^3 P_3(x,y) + t^4\lambda^4 P_4(x,y)\\
&= 6t^2 P_2(x,\mu x + \lambda y) + 4t^3 P_3(x,\mu x + \lambda y) + t^4 P_4(x,\mu x + \lambda y),
\end{aligned}
\]
%CONT>>
